\documentclass[aspectratio=169,10pt]{beamer}

% Shared Beamer setup for Fall 2026 course slides.
% Clean Cal Poly-inspired theme: no custom class files or uncommon packages.
\mode<presentation>{
  \usetheme{default}
  \usefonttheme{professionalfonts}
}

\definecolor{CalPolyGreen}{HTML}{154734}
\definecolor{CalPolyGold}{HTML}{BD8B13}
\definecolor{CalPolySage}{HTML}{7FA28F}
\definecolor{CalPolyMint}{HTML}{DDF1E7}
\definecolor{CalPolySand}{HTML}{A29F86}
\definecolor{CalPolyBeige}{HTML}{EFEEDE}
\definecolor{DeckInk}{HTML}{1F2933}
\definecolor{DeckMuted}{HTML}{5F6B7A}
\definecolor{DeckSoft}{HTML}{F7F7F5}

\setbeamersize{text margin left=0.55in,text margin right=0.55in}
\setbeamercolor{normal text}{fg=DeckInk,bg=white}
\setbeamercolor{structure}{fg=CalPolyGreen}
\setbeamercolor{title}{fg=white}
\setbeamercolor{frametitle}{fg=CalPolyGreen,bg=white}
\setbeamercolor{block title}{bg=CalPolySage,fg=white}
\setbeamercolor{block body}{bg=CalPolyMint,fg=black}
\setbeamercolor{alerted text}{fg=CalPolyGold}

\setbeamerfont{title}{size=\Huge,series=\bfseries}
\setbeamerfont{frametitle}{size=\Large,series=\bfseries}
\setbeamerfont{block title}{size=\normalsize,series=\bfseries}

\setbeamertemplate{navigation symbols}{}
\setbeamertemplate{itemize item}{\color{CalPolyGold}$\blacktriangleright$}
\setbeamertemplate{itemize subitem}{\color{CalPolyGreen}$\bullet$}
\setbeamertemplate{footline}{
  \hfill{\color{DeckMuted}\scriptsize\insertframenumber/\inserttotalframenumber}
  \hspace{0.18in}\vspace{0.12in}
}
\setbeamertemplate{frametitle}{
  \vspace{0.18in}
  \hspace{0.02in}\insertframetitle\par
  \vspace{0.08in}
  \noindent\makebox[\linewidth][l]{%
    {\color{CalPolyGold}\rule{0.55in}{2pt}}%
    {\color{CalPolyGreen}\rule{\dimexpr\linewidth-0.55in\relax}{2pt}}%
  }
}

\newcommand{\CourseNumber}{Course}
\newcommand{\CourseTitle}{Course Title}
\newcommand{\LectureNumber}{0}
\newcommand{\LectureTitle}{Lecture Title}
\newcommand{\LectureDate}{Date}
\newcommand{\CourseUnit}{Unit}

\newcommand{\coursenumber}[1]{\renewcommand{\CourseNumber}{#1}}
\newcommand{\coursetitle}[1]{\renewcommand{\CourseTitle}{#1}}
\newcommand{\lecturenumber}[1]{\renewcommand{\LectureNumber}{#1}}
\newcommand{\lecturetitle}[1]{\renewcommand{\LectureTitle}{#1}}
\newcommand{\lecturedate}[1]{\renewcommand{\LectureDate}{#1}}
\newcommand{\courseunit}[1]{\renewcommand{\CourseUnit}{#1}}

\author{Kirk Duran}
\institute{California Polytechnic State University, San Luis Obispo}
\date{\LectureDate}

\newcommand{\makedecktitle}{
  \begingroup
  \setbeamercolor{background canvas}{bg=CalPolyGreen}
  \begin{frame}[plain]
    \vfill
    \begin{center}
      {\usebeamerfont{title}\color{white}\LectureTitle\par}
      \vspace{0.18in}
      {\Large\color{white}Lecture \LectureNumber\par}
    \end{center}
    \vfill
    \noindent{\color{white}\rule{\linewidth}{1.2pt}}\par
    \vspace{0.08in}
    \noindent\colorbox{CalPolyGold}{\makebox[0.24\linewidth][c]{\color{white}\Large\CourseNumber}}
    \hspace{0.12in}{\color{white}\Large Kirk Duran}
  \end{frame}
  \endgroup
}

\newenvironment{keyidea}{\begin{block}{Key idea}}{\end{block}}
\newenvironment{activity}{\begin{block}{In class}}{\end{block}}
\newenvironment{cpdefinition}{\begin{block}{Definition}}{\end{block}}
\newenvironment{exampleblockcp}{
  \begingroup
  \setbeamercolor{block title}{bg=CalPolySand,fg=white}
  \setbeamercolor{block body}{bg=CalPolyBeige,fg=black}
  \begin{block}{Example}
}{
  \end{block}
  \endgroup
}

\newcommand{\imageplaceholder}[2][0.3in]{%
  \noindent\makebox[\linewidth][c]{%
    \fcolorbox{CalPolySage}{CalPolyBeige}{%
      \parbox[c][#1][c]{0.86\linewidth}{%
        \centering
        {\color{DeckMuted}\scriptsize Image placeholder: }%
        {\color{DeckInk}\scriptsize\ttfamily\detokenize{#2}}%
      }%
    }%
  }%
}

\newcommand{\sectionframe}[1]{
  \begingroup
  \setbeamercolor{background canvas}{bg=CalPolyGreen}
  \begin{frame}[plain]
    \vfill
    {\color{white}\LARGE\bfseries #1\par}
    \vspace{0.18in}
    {\color{CalPolyGold}\rule{0.22\paperwidth}{3pt}}
    {\color{white}\rule{0.62\paperwidth}{3pt}}
    \vfill
  \end{frame}
  \endgroup
}

\newcommand{\placeholdernote}{
  \begin{block}{Placeholder}
    Replace these bullets with the final lecture material, examples, and in-class activities.
  \end{block}
}

\usepackage[T1]{fontenc}
\usepackage{lmodern}
\pdfstringdefDisableCommands{\def\translate#1{#1}}

\graphicspath{{../assets/lecture-26/}}

% Generated teaching diagrams; prompts are recorded alongside the assets.
\newcommand{\lectureimage}[2][1.75in]{%
  \centering
  \includegraphics[width=\linewidth,height=#1,keepaspectratio]{#2}%
}

\setbeamersize{text margin left=0.42in,text margin right=0.42in}

\coursenumber{CS 1001}
\coursetitle{Introduction to Computing}
\lecturenumber{26}
\lecturetitle{Multidimensional Data}
\lecturedate{Fall 2026}
\courseunit{10. Structured Data}

\begin{document}

\makedecktitle

\begin{frame}{Today}
  \begin{itemize}
    \item Extend the sequence model from one dimension to data organized by multiple coordinates.
    \item Represent two-dimensional data using nested Python lists.
    \item Read and mutate elements using row and column indexes.
    \item Trace nested iteration over rows and individual elements.
    \item Generalize the same model to three and more dimensions.
    \item Introduce NumPy arrays as a regular multidimensional representation for numerical data.
  \end{itemize}

  \begin{keyidea}
    Multidimensional data does not require a new notion of indexing. It requires applying familiar indexing and iteration at more than one structural level.
  \end{keyidea}
\end{frame}

\begin{frame}{The Boundary for This Lecture}
  \begin{columns}[T,onlytextwidth]
    \begin{column}{0.47\textwidth}
      \begin{block}{Today}
        \begin{itemize}
          \item nested lists
          \item rows, columns, and coordinates
          \item two-stage indexing
          \item mutation and aliasing in nested lists
          \item nested iteration
          \item 3D and $n$-dimensional structure
          \item NumPy arrays, shape, slicing, and reshape
        \end{itemize}
      \end{block}
    \end{column}
    \begin{column}{0.47\textwidth}
      \begin{alertblock}{Not today}
        \begin{itemize}
          \item general-purpose built-in function catalogs
          \item searching or sorting algorithms
          \item algorithmic complexity
          \item advanced NumPy broadcasting rules
          \item functions over structured data
          \item object-oriented representations
        \end{itemize}
      \end{alertblock}
    \end{column}
  \end{columns}
\end{frame}

\sectionframe{Part 1: From One Dimension to Multiple Dimensions}

\begin{frame}{A One-Dimensional Sequence Uses One Position}
  \begin{columns}[T,onlytextwidth]
    \begin{column}{0.54\textwidth}
      \begin{block}{One-dimensional data}
        \centering
        \texttt{scores = [88, 94, 79, 91]}
      \end{block}

      \begin{itemize}
        \item One integer index identifies a position.
        \item \texttt{scores[2]} selects one element.
        \item The structure has one natural positional axis: left to right.
      \end{itemize}
    \end{column}
    \begin{column}{0.40\textwidth}
      % Intended visual: one horizontal sequence labeled with positions 0,1,2,3.
      \lectureimage[2.0in]{one-dimensional-sequence-axis.png}
    \end{column}
  \end{columns}
\end{frame}

\begin{frame}{Two Coordinates Identify a Cell}
  Seats, pixels, and board squares are naturally identified by two coordinates.

  \begin{center}
    \renewcommand{\arraystretch}{1.28}
    \begin{tabular}{c|ccc}
      & \textbf{Column 0} & \textbf{Column 1} & \textbf{Column 2} \\\hline
      \textbf{Row 0} & 4 & 7 & 1 \\
      \textbf{Row 1} & 8 & 5 & 2 \\
      \textbf{Row 2} & 6 & 3 & 9
    \end{tabular}
  \end{center}

  \begin{cpdefinition}
    \textbf{Multidimensional data} is data whose organization requires more than one positional coordinate to identify a particular component.
  \end{cpdefinition}

  \begin{keyidea}
    Two coordinates identify a cell: row first, then column.
  \end{keyidea}
\end{frame}

\begin{frame}{A Python 2D List Is a List Whose Elements Are Lists}
  \begin{block}{Nested list}
    \texttt{grid = [[4, 7, 1],}\\
    \hspace*{0.78in}\texttt{[8, 5, 2],}\\
    \hspace*{0.78in}\texttt{[6, 3, 9]]}
  \end{block}

  \begin{itemize}
    \item The outer list contains three elements.
    \item Each outer element is itself a list.
    \item We interpret each inner list as one row.
    \item This creates a useful two-dimensional \emph{conceptual} structure from ordinary lists.
  \end{itemize}

  \begin{keyidea}
    Python does not introduce a special ``2D list'' type. The representation is nested ordinary lists.
  \end{keyidea}
\end{frame}

\begin{frame}{Rectangular and Ragged Nested Lists}
  \begin{columns}[T,onlytextwidth]
    \begin{column}{0.47\textwidth}
      \begin{block}{Rectangular}
        \texttt{[[1, 2, 3],}\\
        \hspace*{0.16in}\texttt{[4, 5, 6],}\\
        \hspace*{0.16in}\texttt{[7, 8, 9]]}
      \end{block}
      Every row has the same number of elements.
    \end{column}
    \begin{column}{0.47\textwidth}
      \begin{block}{Ragged}
        \texttt{[[1, 2],}\\
        \hspace*{0.16in}\texttt{[3, 4, 5],}\\
        \hspace*{0.16in}\texttt{[6]]}
      \end{block}
      Different rows may have different lengths.
    \end{column}
  \end{columns}

  \begin{alertblock}{Consequence}
    A nested list can represent a grid, but Python does not require it to be rectangular.
  \end{alertblock}
\end{frame}

\begin{frame}{The Row--Column Coordinate Model}
  \begin{columns}[T,onlytextwidth]
    \begin{column}{0.52\textwidth}
      For a rectangular nested list, we use the convention:

      \begin{center}
        \Large\texttt{grid[row][column]}
      \end{center}

      \begin{itemize}
        \item first index: which row?
        \item second index: which element inside that row?
        \item both indexes remain zero-based
      \end{itemize}
    \end{column}
    \begin{column}{0.42\textwidth}
      % Intended visual: 3x3 grid with row 1 / column 2 highlighted and coordinate arrows.
      \lectureimage[2.1in]{row-column-coordinate-grid.png}
    \end{column}
  \end{columns}
\end{frame}

\sectionframe{Part 2: Indexing, Mutation, and References}

\begin{frame}{Two-Dimensional Indexing Is Two Ordinary Indexing Operations}
  Suppose:

  \begin{block}{Data}
    \texttt{grid = [[4, 7, 1], [8, 5, 2], [6, 3, 9]]}
  \end{block}

  Evaluate:

  \begin{center}
    \Large\texttt{grid[1][2]}
  \end{center}

  \begin{enumerate}
    \item \texttt{grid[1]} evaluates to the row \texttt{[8, 5, 2]}.
    \item The remaining \texttt{[2]} indexes that row.
    \item The final value is \texttt{2}.
  \end{enumerate}
\end{frame}

\begin{frame}{Trace a Two-Index Expression}
  \begin{columns}[T,onlytextwidth]
    \begin{column}{0.52\textwidth}
      \begin{block}{Expression}
        \centering
        \texttt{grid[2][1]}
      \end{block}

      \begin{center}
        \texttt{grid[2]} $\rightarrow$ \texttt{[6, 3, 9]}\\[0.12in]
        \texttt{[6, 3, 9][1]} $\rightarrow$ \texttt{3}
      \end{center}

      \begin{keyidea}
        Bracket operations associate left to right: obtain the inner list first, then index that list.
      \end{keyidea}
    \end{column}
    \begin{column}{0.42\textwidth}
      % Intended visual: outer list arrow to row list, then second arrow to one element.
      \lectureimage[2.1in]{nested-list-two-stage-indexing.png}
    \end{column}
  \end{columns}
\end{frame}

\begin{frame}{One Index Selects an Entire Row}
  \begin{block}{Example}
    \texttt{grid = [[4, 7, 1], [8, 5, 2], [6, 3, 9]]}\\[0.10in]
    \texttt{row = grid[1]}
  \end{block}

  Afterwards:

  \begin{center}
    \texttt{row} $\rightarrow$ \texttt{[8, 5, 2]}
  \end{center}

  \begin{itemize}
    \item The outer list stores references to its element values.
    \item Here, the selected element value is itself a list.
    \item \texttt{row} and \texttt{grid[1]} therefore refer to the same inner list.
  \end{itemize}
\end{frame}

\begin{frame}{Mutating One Cell}
  \begin{block}{Before}
    \texttt{grid = [[4, 7, 1], [8, 5, 2], [6, 3, 9]]}
  \end{block}

  \begin{block}{Mutation}
    \texttt{grid[1][2] = 20}
  \end{block}

  \begin{block}{After}
    \texttt{[[4, 7, 1], [8, 5, 20], [6, 3, 9]]}
  \end{block}

  \begin{keyidea}
    Indexed assignment mutates the selected inner list. The surrounding outer list still refers to that same row list.
  \end{keyidea}
\end{frame}

\begin{frame}{The Reference Model Becomes More Important When Lists Are Nested}
  \begin{columns}[T,onlytextwidth]
    \begin{column}{0.52\textwidth}
      Conceptually:

      \begin{center}
        \texttt{grid} $\longrightarrow$ outer list\\[0.08in]
        outer positions $\longrightarrow$ row lists\\[0.08in]
        row positions $\longrightarrow$ element values
      \end{center}

      \begin{itemize}
        \item The rows are separate list values only if they were created separately.
        \item Multiple outer positions can refer to the same row list.
      \end{itemize}
    \end{column}
    \begin{column}{0.42\textwidth}
      % Intended visual: outer list containing three arrows to three distinct row lists.
      \lectureimage[2.1in]{nested-list-reference-model.png}
    \end{column}
  \end{columns}
\end{frame}

\begin{frame}{Aliasing Can Occur Inside a Nested List}
  \begin{block}{One row, referenced three times}
    \texttt{row = [0, 0, 0]}\\
    \texttt{grid = [row, row, row]}\\[0.08in]
    \texttt{grid[0][1] = 7}
  \end{block}

  What is \texttt{grid} afterward?

  \begin{center}
    \Large\texttt{[[0, 7, 0], [0, 7, 0], [0, 7, 0]]}
  \end{center}

  \begin{alertblock}{Why?}
    The outer list contains three references to the same inner list. The mutation changes that one shared row.
  \end{alertblock}

  To make rows independent, create separate inner lists, for example \texttt{[[0, 0, 0], [0, 0, 0], [0, 0, 0]]}.
\end{frame}

\begin{frame}{Row Slicing Uses Ordinary List Slicing}
  \begin{block}{Data}
    \texttt{grid = [[1, 2, 3], [4, 5, 6], [7, 8, 9]]}
  \end{block}

  \begin{block}{Select rows 1 through the end}
    \texttt{rows = grid[1:]}
  \end{block}

  The new outer list contains references to:

  \begin{center}
    \texttt{[[4, 5, 6], [7, 8, 9]]}
  \end{center}

  \begin{alertblock}{Reference detail}
    Slicing creates a new outer list, but the selected inner row lists are still shared with \texttt{grid}.
  \end{alertblock}

  Native nested lists have no special column slice such as \texttt{grid[:, 1]}; a column must be selected from the rows during iteration. NumPy will provide direct multidimensional slicing later.
\end{frame}

\sectionframe{Part 3: Nested Iteration}

\begin{frame}{The Outer Loop Naturally Visits Rows}
  \begin{block}{Code}
    \texttt{grid = [[1, 2, 3], [4, 5, 6], [7, 8, 9]]}\\[0.08in]
    \texttt{for row in grid:}\\
    \hspace*{0.20in}\texttt{...}
  \end{block}

  Successive bindings of \texttt{row} are:

  \begin{center}
    \texttt{[1, 2, 3]} $\rightarrow$ \texttt{[4, 5, 6]} $\rightarrow$ \texttt{[7, 8, 9]}
  \end{center}

  \begin{keyidea}
    Iterating over the outer list visits its elements. Because those elements are row lists, each loop iteration receives one row.
  \end{keyidea}
\end{frame}

\begin{frame}{A Nested Loop Visits Individual Cells}
  \begin{block}{Code}
    \texttt{for row in grid:}\\
    \hspace*{0.20in}\texttt{for value in row:}\\
    \hspace*{0.40in}\texttt{...}
  \end{block}

  \begin{center}
    outer loop $\rightarrow$ selects one row\\[0.08in]
    inner loop $\rightarrow$ visits each value in that row
  \end{center}

  \begin{keyidea}
    The inner loop completes for the current row before the outer loop advances to the next row.
  \end{keyidea}
\end{frame}

\begin{frame}{Trace Nested Iteration}
  For \texttt{grid = [[1, 2], [3, 4]]}:

  \begin{center}
    \renewcommand{\arraystretch}{1.30}
    \begin{tabular}{c|c|c}
      \textbf{Step} & \textbf{row} & \textbf{value} \\\hline
      1 & \texttt{[1, 2]} & \texttt{1} \\
      2 & \texttt{[1, 2]} & \texttt{2} \\
      3 & \texttt{[3, 4]} & \texttt{3} \\
      4 & \texttt{[3, 4]} & \texttt{4}
    \end{tabular}
  \end{center}

  \begin{block}{Control-flow fact}
    The outer binding remains the same while the inner loop advances through the current row.
  \end{block}
\end{frame}

\begin{frame}{Index-Based Nested Iteration Exposes Coordinates}
  For a known $3 \times 3$ grid:

  \begin{block}{Code}
    \texttt{for row in range(3):}\\
    \hspace*{0.20in}\texttt{for col in range(3):}\\
    \hspace*{0.40in}\texttt{value = grid[row][col]}
  \end{block}

  \begin{itemize}
    \item \texttt{row} identifies the outer-list position.
    \item \texttt{col} identifies the position inside that row.
    \item The pair \texttt{(row, col)} identifies one cell.
  \end{itemize}

  \begin{keyidea}
    Direct iteration emphasizes values; index-based iteration is useful when the coordinates themselves matter.
  \end{keyidea}
\end{frame}

\begin{frame}[shrink=4]{Example Representation: A Tic-Tac-Toe Board}
  \begin{columns}[T,onlytextwidth]
    \begin{column}{0.52\textwidth}
      \begin{block}{Board}
        \texttt{board = [["X", "", "O"],}\\
        \hspace*{0.54in}\texttt{["", "X", ""],}\\
        \hspace*{0.54in}\texttt{["O", "", "X"]]}
      \end{block}

      \begin{itemize}
        \item row and column coordinates identify squares;
        \item an empty string can represent an unoccupied square;
        \item mutation can represent a move.
      \end{itemize}
    \end{column}
    \begin{column}{0.42\textwidth}
      % Intended visual: tic-tac-toe board annotated with row/column coordinates.
      \lectureimage[1.8in]{tic-tac-toe-coordinate-board.png}
    \end{column}
  \end{columns}

  \begin{keyidea}
    The data representation is the topic here; game-winning logic and algorithms are outside today's scope.
  \end{keyidea}
\end{frame}

\sectionframe{Part 4: Three and More Dimensions}

\begin{frame}{A Third Dimension Adds Another Structural Level}
  A three-dimensional nested list can be interpreted as a sequence of two-dimensional layers.

  \begin{center}
    \textbf{layer} $\rightarrow$ \textbf{row} $\rightarrow$ \textbf{column} $\rightarrow$ \textbf{value}
  \end{center}

  Suppose:

  \begin{block}{Data}
    \texttt{cube = [[[1, 2], [3, 4]],}\\
    \hspace*{0.64in}\texttt{[[5, 6], [7, 8]]]}
  \end{block}

  Evaluate \texttt{cube[1][0][1]}:

  \begin{enumerate}
    \item \texttt{cube[1]} $\rightarrow$ \texttt{[[5, 6], [7, 8]]}
    \item \texttt{[0]} $\rightarrow$ \texttt{[5, 6]}
    \item \texttt{[1]} $\rightarrow$ \texttt{6}
  \end{enumerate}

  \begin{keyidea}
    Additional dimensions repeat the same rule: each index moves one structural level deeper.
  \end{keyidea}
\end{frame}

\begin{frame}[shrink=4]{Nested Iteration Mirrors Structural Depth}
  \begin{columns}[T,onlytextwidth]
    \begin{column}{0.55\textwidth}
      \begin{block}{Three levels}
        \texttt{for layer in cube:}\\
        \hspace*{0.20in}\texttt{for row in layer:}\\
        \hspace*{0.40in}\texttt{for value in row:}\\
        \hspace*{0.60in}\texttt{...}
      \end{block}

      Each loop operates on the sequence value produced by the level above it.
    \end{column}
    \begin{column}{0.39\textwidth}
      % Intended visual: layer -> row -> cell hierarchy with three nested boxes.
      \lectureimage[1.8in]{three-dimensional-layer-row-column.png}
    \end{column}
  \end{columns}

  \begin{keyidea}
    Each additional nested sequence level can require another coordinate and another level of iteration. This same reasoning generalizes beyond three dimensions.
  \end{keyidea}
\end{frame}

\begin{frame}{Nested Lists Are Flexible---Sometimes Too Flexible}
  Python nested lists permit:

  \begin{itemize}
    \item rows of different lengths;
    \item inner lists that are shared through aliases;
    \item elements of different types;
    \item arbitrary nesting depth.
  \end{itemize}

  That flexibility is useful for general programming, but regular numerical data often benefits from a structure with explicit dimensional shape.

  \begin{keyidea}
    This is the point where NumPy provides a different abstraction: the multidimensional array.
  \end{keyidea}
\end{frame}

\sectionframe{Part 5: NumPy Multidimensional Arrays}

\begin{frame}{Why Introduce NumPy Here?}
  \begin{columns}[T,onlytextwidth]
    \begin{column}{0.55\textwidth}
      NumPy is designed around regular numerical arrays.

      \begin{itemize}
        \item dimensional shape is explicit;
        \item elements use a common numerical data type;
        \item multidimensional indexing is built into the array abstraction;
        \item arithmetic can operate over many array elements at once.
      \end{itemize}
    \end{column}
    \begin{column}{0.39\textwidth}
      % Intended visual: irregular nested lists contrasted with a regular rectangular ndarray.
      \lectureimage[2.1in]{nested-lists-versus-numpy-array.png}
    \end{column}
  \end{columns}
\end{frame}

\begin{frame}{Creating a NumPy Array}
  \begin{block}{Import and construction}
    \texttt{import numpy as np}\\[0.10in]
    \texttt{matrix = np.array([[1, 2, 3],}\\
    \hspace*{1.22in}\texttt{[4, 5, 6]])}
  \end{block}

  \begin{itemize}
    \item \texttt{np.array(...)} constructs a NumPy array from the provided data.
    \item The result is a two-dimensional NumPy \texttt{ndarray}.
    \item The rows and columns are part of the array's explicit shape.
  \end{itemize}

  \begin{alertblock}{Scope}
    We use NumPy only to introduce multidimensional array operations needed for this unit; this is not a full NumPy API survey.
  \end{alertblock}
\end{frame}

\begin{frame}{Shape Describes the Size of Each Dimension}
  For:

  \begin{block}{Array}
    \texttt{matrix = np.array([[1, 2, 3], [4, 5, 6]])}
  \end{block}

  \begin{center}
    \texttt{matrix.shape} $\rightarrow$ \texttt{(2, 3)}
  \end{center}

  Interpret the tuple as:

  \begin{center}
    \textbf{2 rows} $\times$ \textbf{3 columns}
  \end{center}

  \begin{keyidea}
    Shape makes the dimensional structure explicit rather than leaving it as an informal interpretation of nested lists.
  \end{keyidea}
\end{frame}

\begin{frame}{NumPy Exposes Dimension Count and Element Type}
  \begin{block}{Given}
    \texttt{matrix = np.array([[1, 2, 3], [4, 5, 6]])}
  \end{block}

  \begin{center}
    \renewcommand{\arraystretch}{1.35}
    \begin{tabular}{l|l}
      \textbf{Attribute} & \textbf{Meaning} \\\hline
      \texttt{matrix.ndim} & number of dimensions \\
      \texttt{matrix.shape} & size along each dimension \\
      \texttt{matrix.dtype} & stored element data type
    \end{tabular}
  \end{center}

  \begin{block}{For this array}
    \texttt{ndim} is \texttt{2}; \texttt{shape} is \texttt{(2, 3)}.
  \end{block}
\end{frame}

\begin{frame}{NumPy Uses Multidimensional Index Syntax}
  \begin{columns}[T,onlytextwidth]
    \begin{column}{0.47\textwidth}
      \begin{block}{Nested list}
        \centering
        \texttt{grid[1][2]}
      \end{block}
      Two ordinary indexing operations.
    \end{column}
    \begin{column}{0.47\textwidth}
      \begin{block}{NumPy array}
        \centering
        \texttt{matrix[1, 2]}
      \end{block}
      One array indexing operation with two coordinates.
    \end{column}
  \end{columns}

  \begin{center}
    row coordinate first, column coordinate second
  \end{center}
\end{frame}

\begin{frame}{NumPy Can Slice Along Individual Dimensions}
  Suppose \texttt{matrix} has shape \texttt{(3, 4)}.

  \begin{center}
    \renewcommand{\arraystretch}{1.35}
    \begin{tabular}{l|l}
      \textbf{Expression} & \textbf{Selection} \\\hline
      \texttt{matrix[1, :]} & all columns from row 1 \\
      \texttt{matrix[:, 2]} & column 2 from all rows \\
      \texttt{matrix[0:2, 1:3]} & a rectangular subregion
    \end{tabular}
  \end{center}

  \begin{keyidea}
    Each position in the comma-separated index controls one array dimension.
  \end{keyidea}
\end{frame}

\begin{frame}{Visualizing Row and Column Slices}
  \begin{columns}[T,onlytextwidth]
    \begin{column}{0.42\textwidth}
      \begin{block}{Row slice}
        \centering
        \texttt{matrix[1, :]}
      \end{block}

      \begin{block}{Column slice}
        \centering
        \texttt{matrix[:, 2]}
      \end{block}
    \end{column}
    \begin{column}{0.52\textwidth}
      % Intended visual: same matrix shown twice, once with a full row highlighted and once with a full column highlighted.
      \lectureimage[2.1in]{numpy-slicing-rows-columns.png}
    \end{column}
  \end{columns}
\end{frame}

\begin{frame}{Array Arithmetic Can Express Elementwise Transformation}
  \begin{block}{Array}
    \texttt{matrix = np.array([[1, 2], [3, 4]])}
  \end{block}

  \begin{block}{Elementwise operation}
    \texttt{doubled = matrix * 2}
  \end{block}

  \begin{center}
    \texttt{doubled} $\rightarrow$
    \texttt{[[2, 4], [6, 8]]}
  \end{center}

  \begin{keyidea}
    NumPy lets an operation describe a transformation over the array as a whole instead of requiring an explicit nested Python loop for this simple case.
  \end{keyidea}
\end{frame}

\begin{frame}{Reshape Changes Dimensional Organization}
  \begin{columns}[T,onlytextwidth]
    \begin{column}{0.56\textwidth}
      \begin{block}{Start with six values}
        \texttt{data = np.array([1, 2, 3,}\\
        \hspace*{0.30in}\texttt{4, 5, 6])}
      \end{block}

      \begin{block}{Reshape}
        \texttt{matrix = data.reshape((2, 3))}
      \end{block}

      \begin{keyidea}
        The new shape must preserve the number of elements: $6=2\times3$.
      \end{keyidea}
    \end{column}
    \begin{column}{0.38\textwidth}
      % Intended visual: six cells in one row rearranged into a 2x3 rectangle.
      \lectureimage[2.0in]{reshape-same-data-new-shape.png}
    \end{column}
  \end{columns}
\end{frame}

\begin{frame}{Matrix Multiplication Is Different from Elementwise Multiplication}
  For compatible two-dimensional arrays:

  \begin{block}{Matrix product}
    \centering
    \texttt{result = a @ b}
  \end{block}

  \begin{itemize}
    \item \texttt{a * b} means elementwise multiplication in NumPy.
    \item \texttt{a @ b} means the linear-algebra matrix product.
    \item Matrix dimensions must be compatible for the product to be defined.
  \end{itemize}

  \begin{alertblock}{Scope}
    We use this only to distinguish two common array operations. Matrix-multiplication algorithms are outside this course unit.
  \end{alertblock}
\end{frame}

\begin{frame}{Nested Lists and NumPy Arrays Serve Different Roles}
  \begin{center}
    \renewcommand{\arraystretch}{1.30}
    \begin{tabular}{l|l|l}
      \textbf{Property} & \textbf{Nested lists} & \textbf{NumPy ndarray} \\\hline
      Structure & flexible nesting & explicit dimensions \\
      Rectangular required? & no & regular shape \\
      Element types & may vary & common dtype \\
      2D indexing & \texttt{a[r][c]} & \texttt{a[r, c]} \\
      Column slicing & manual iteration & direct multidimensional slice \\
      Numeric whole-array ops & explicit iteration & array operations
    \end{tabular}
  \end{center}

  \begin{keyidea}
    The choice of representation should match the structure and operations the program needs.
  \end{keyidea}
\end{frame}

\begin{frame}{Integrated Prediction}
  Consider:

  \begin{block}{Code}
    \texttt{matrix = np.array([[2, 4, 6],}\\
    \hspace*{1.22in}\texttt{[8, 10, 12]])}\\[0.08in]
    \texttt{column = matrix[:, 1]}\\
    \texttt{changed = matrix * 2}
  \end{block}

  Predict:
  \begin{enumerate}
    \item \texttt{matrix.shape}
    \item the value selected by \texttt{matrix[1, 2]}
    \item the values in \texttt{column}
    \item the values in \texttt{changed}
    \item whether \texttt{matrix * 2} changes the original \texttt{matrix}
  \end{enumerate}
\end{frame}

\begin{frame}{Multidimensional Data: The Complete Picture}
  \begin{center}
    \renewcommand{\arraystretch}{1.32}
    \begin{tabular}{l|l}
      \textbf{Representation} & \textbf{Core reasoning model} \\\hline
      1D list & one index selects one element \\
      nested 2D list & first index selects row, second selects cell \\
      nested 3D list & layer $\rightarrow$ row $\rightarrow$ cell \\
      NumPy array & one coordinate per explicit array dimension
    \end{tabular}
  \end{center}

  \begin{keyidea}
    More dimensions add coordinates, not a fundamentally different execution model. Trace one structural level at a time.
  \end{keyidea}
\end{frame}

\begin{frame}{Takeaways}
  \begin{enumerate}
    \item A two-dimensional Python list is a nested list: an outer list whose elements are row lists.
    \item \texttt{grid[row][col]} is evaluated as two ordinary indexing operations.
    \item Nested mutable structures require careful reference reasoning because inner lists can be aliased.
    \item Nested loops mirror nested structure: each loop descends one sequence level.
    \item Three and more dimensions extend the same coordinate model.
    \item NumPy arrays make regular multidimensional shape explicit and support multidimensional indexing, slicing, arithmetic, and reshaping.
  \end{enumerate}

  \begin{block}{Next}
    Lecture 27 asks how to design and reason about \textbf{functions over structured data}.
  \end{block}
\end{frame}

\end{document}
